\section{Valuation and Profit and Loss}\label{ch:valuation}

This chapter discharges §8 of the constitution, as amended by Amendment~1.

Net Asset Value (NAV) has an operational meaning before it has a formula.
The initial NAV of a wallet is the cash it took to build the portfolio at
inception; the current NAV is the cash you would hold if you unwound every
position at prevailing prices; and the profit and loss (PnL) is the
difference between the two. Nothing in this chapter is a new store: every
number it produces is a projection of the immutable log, computed on demand
and never written back.

\subsection{Net Asset Value}

The operational account becomes a formula in two steps. First, replay the log
to recover the wallet's current composition --- the signed coordinate vector,
per unit, that the fold of Chapter~2 produces. Second, multiply that
composition by the current price vector supplied by the pricing data layer,
summing over the valued units only.

\begin{definition}[Net Asset Value]
The NAV of a wallet $w$ at a log position is
\[
  \mathrm{NAV}(w) \;=\; \sum_{u \,\text{valued}} \mathrm{owned}(w,u)\,\cdot\, P\!\left(u, \sigma(u)\right),
\]
where $\mathrm{owned}(w,u)$ is the owned coordinate of $w$'s balance in unit
$u$ (Chapter~3), $\sigma(u)$ is the current lifecycle state of $u$ read from
the ledger, and $P$ is the price the pricing data layer supplies for that unit
in that state.
\end{definition}

Only the owned coordinate enters. The lent and posted coordinates carry mass,
never value (Chapter~3): a security posted as collateral is still owned by the
poster and is already counted once, on its owned coordinate; counting the
posted marker as well would value one economic reality twice. Non-valued units
--- a strategy rulebook, a mandate's governance shell --- have no price and
enter no sum; their value is undefined, not zero.

Two disciplines are visible in the definition and are carried by the machinery
of Chapters~3 and~8. The composition is one number per coordinate, obtained
from one fold of one log, so no two readers can disagree about what is held.
And a price is never combined with a quantity across a corporate-action event:
after a two-for-one split, $2{,}000$ shares are multiplied by the post-split
frame $50.00$, never by the stale $100.00$ --- consistency of reference,
catalogued in Chapter~14.

\begin{lstlisting}
type Points = Integer                      -- money in minor units, exact
price :: Unit -> State -> PriceVector -> Points   -- pricing is state-parameterised
nav   :: Composition -> PriceVector -> Points
nav c pv = sum [ q * price u (stateOf c u) pv     -- owned coordinate, valued units
               | (u, q) <- ownedBalances c, valued u ]
\end{lstlisting}

\subsection{Profit and loss is the variation in NAV}

\begin{principle}[PnL is NAV variation]
The profit and loss of any wallet between two log positions is exactly the
variation in its NAV over that interval. PnL is therefore path-independent,
and no second copy of ``realised PnL'' or ``book cost'' is ever maintained
alongside the position record.
\end{principle}

The principle is a consequence of the definition, not an additional
convention. NAV is a function of the composition and the price vector alone;
its variation between two positions depends only on the endpoints, never on
the sequence of transactions that connected them. A stored realised-PnL figure
would be a second record of a fact the log already determines --- precisely
the redundancy the framework exists to remove. Realised and unrealised PnL
remain available as \emph{projections}, cut from the same log by a reporting
convention (Chapter~12), but neither is a source of truth.

\paragraph{Episode --- the future.}
FUT-IDX is a cash-settled future on index IDX, multiplier $10$, USD,
settled-to-market. Wallet W-ALPHA buys one contract at $995$ on day~0. The
official settlement prints on the three settlement days are $1{,}000$, then
$990$, then $1{,}005$; the contract's smart contract fires on each print and
emits the day's variation-margin (VM) move.

\begin{center}
\begin{tabular}{lrrr}
\toprule
Settlement day & Print & VM move (owned cash) & Running NAV change \\
\midrule
Day 1 & $1{,}000$ & $+50$ & $+50$ \\
Day 2 & $\phantom{0}990$ & $-100$ & $-50$ \\
Day 3 & $1{,}005$ & $+150$ & $+100$ \\
\bottomrule
\end{tabular}
\end{center}

Each VM is a settlement --- an owned-cash move that extinguishes the day's
exposure, with no return obligation, because none exists in law for a
settled-to-market contract (§8, case~1; developed in Chapter~9). The
cumulative VM is $50 - 100 + 150 = 100$, and the position's PnL is
$(1{,}005 - 995)\times 10 = 100$. The two agree, and the second computation
never touches the intermediate prints $1{,}000$ and $990$: the path is
irrelevant, only the endpoints $995$ and $1{,}005$ matter.
\emph{PnL is the variation in NAV, and it is path-independent.}

\subsection{Pricing is parameterised by state read from the ledger}

Because a unit's value depends on where it stands in its lifecycle, the price
function is parameterised not only by the unit but by its state --- and that
state is obtained from the ledger, from UnitStatus and PositionState
(Chapter~6), never from a store the pricing layer keeps of its own. The pricing
layer computes; the ledger is the sole authority for the state it computes in.

\paragraph{Episode --- the one-touch.}
OT-1 is a one-touch on stock OMEGA, barrier $80.00$ observed at the official
close (a discretely-monitored closing-price one-touch), payout
\USD{1,000,000}, held by W-ALPHA. While live, the pricing layer marks it at
\USD{400,000}. When OMEGA's official close prints $79.00$, the Event Monitor
emits the touch and the option's contract records the lifecycle transition:
UnitStatus $\rightarrow$ \emph{triggered}. The pricing layer, reading that
state from the ledger, now prices the same unit in state \emph{triggered} at
$0$: the contingent claim has resolved, its economic value realised as the
cash payout that Chapter~9 discharges.

The mark on OT-1 collapses from \USD{400,000} to $0$, but the mass does not
move: one unit of OT-1 remains on W-ALPHA's owned coordinate, and any posted
marker over it remains on its posted coordinate, until the payout completes and
the unit retires from the zero vector. \emph{A lifecycle event extinguishes
value, never mass: coordinates carry mass, prices carry value, and the price
turns on a state that lives on the record.}

\subsection{Deposits do not move profit and loss}

For the variation in NAV to be exactly profit, an inflow of value must not, by
itself, create profit. The constitution's guarantee (§8, as amended) is that
every inflow of value is one of three things, and which one is not a matter of
asset class but of the legal regime declared once on the governing agreement
unit.

\begin{enumerate}[label=(\arabic*),leftmargin=*]
\item \textbf{An exchange paired with an equal-valued outflow.}
  Settled-to-market margin is this case: an owned-cash move that settles the
  day's exposure with no return obligation, because the settlement
  extinguishes the exposure outright. The future's VM above is the worked
  instance.
\item \textbf{A contribution or financing received against an equal obligation
  created in the same transaction.} Collateral received under title transfer,
  cash included, is this case: the receiver owns the inflow and owes an
  equivalent-return obligation that is itself a unit.
\item \textbf{Value held in custody without being owned.} Collateral received
  under a security interest without right of use, segregated cash included, is
  this case: it is recorded on the received-collateral coordinate and never
  touches owned.
\end{enumerate}

Case~1 adds equal and opposite owned value; case~3 adds no owned value at all.
Case~2 is the one that could leak profit, and does not, for a stated reason.

\begin{proposition}[Deposit-neutrality of title-transfer financing]
The return obligation minted in case~2 is a valued unit, priced at the inflow
amount --- par plus accrued at the declared rate. This is a theorem, and not an
approximation, exactly when the declared financing rate equals the fair rate for
the term, so that par-plus-accrued is the obligation's fair value; a declared rate
away from the fair rate prices the obligation off the inflow amount, and the
day-one difference is a financing basis, not a deposit. Under that assumption, at
receipt the owned cash and the obligation are equal and opposite in the owned sum,
so net owned value is unchanged, and NAV does not move.
\end{proposition}

Because the obligation is priced at par-plus-accrued, its value tracks the cash
it returns for the life of the financing, and the receipt neither raises nor
lowers NAV; only the subsequent accrual --- a genuine financing cost or
benefit --- moves it, as it should. This is why PnL needs no regime branch. The
regime decides \emph{where} an inflow lands --- owned plus an obligation unit,
a posted marker, or an owned settlement --- but once it has landed, PnL is
computed uniformly.

\begin{principle}[No regime branch in PnL]
PnL is $\sum \mathrm{owned}\cdot P$ over the valued units, with no branch on
collateral regime. The regime is a property of the transaction shape and the
agreement unit's declared terms, not of the valuation.
\end{principle}

A deposit therefore cannot move profit and loss --- not because a formula
subtracts it, but because, under every case, it cannot change net owned value.
The full mechanics of the three regimes, the obligation and claim units, and
the entitlement trap are Chapter~9's; this chapter uses only their valuation
consequence.

\subsection{Mid-life valuation is read from the record}

State-parameterised pricing extends past a binary knock to any unit whose value
depends on facts its own history has accrued. The point is that those facts
have a home on the record (Chapter~6), so the pricing layer reads them; it does
not remember them.

\paragraph{Episode --- the variance swap.}
VS-1 is a one-year OTC variance swap on IDX, $252$ scheduled fixings against
the recorded official closes, variance strike $K = 400$ points, variance
notional $N = \USD{1,000}$ per point, no cap. Its declared convention
annualises over $252$ returns, so realised variance in points is
$\sigma^2_R = 10^4 \cdot \sum_{i=1}^{252} R_i^2$, where $R_i$ is the $i$-th log
return of the close series. Each firing writes the running
sum-of-squared-returns statistic into UnitStatus as an integer at scale
$10^{-8}$, so firing $k$ finds fixings $1,\dots,k-1$ already on the record.

VS-1 fixes on the same recorded IDX closes the future settles against. The
future's three settlement days are fixings $124$, $125$, $126$ of the series
(closes $1{,}000$, $990$, $1{,}005$); the future's final settlement level,
$1{,}005$, is therefore the variance swap's mid-life checkpoint close at fixing
$126$ --- one recorded fact, two consumers.

At the mid-life checkpoint (fixing $126$, half the fixings elapsed) the accrued
statistic is $A_{126} = \sum_{i=1}^{126} R_i^2 = 0.018050$, stored as the
integer $1{,}805{,}000$. Expressed in points and annualised to date this is
$10^4 \cdot A_{126}\cdot(252/126) = 361$ points --- a realised $19$-vol over the
first half. The mid-life mark to the variance receiver is that accrued
statistic plus the remaining implied variance, both taken from the record: with
remaining implied variance $I = 484$ points ($22$-vol) over the remaining
half-year fraction $126/252$,
\[
  V_{126} \;=\; N\left[\,10^4\!\cdot A_{126} \;+\; I\cdot\tfrac{126}{252} \;-\; K\,\right]
         \;=\; 1{,}000\left[\,180.5 + 242 - 400\,\right] \;=\; \USD{22,500}.
\]
The accrued term ($180.5$ points) is a projection of a UnitStatus integer; the
implied term ($242$ points) is a market observation with provenance
(Chapter~8); the strike is a declared term. No figure is stored as a mark. When
the year completes, the remaining fixings accrue $10^4\!\cdot\!\sum_{127}^{252}
R_i^2 = 260.5$ points, the realised variance is $180.5 + 260.5 = 441$ points,
and the settlement is $1{,}000\times(441 - 400) = \USD{41,000}$ to the receiver
--- the mid-life mark and the endpoint cut from one accruing statistic.
\emph{A mid-life value is not a stored quantity; it is the record, priced.}

\paragraph{Episode --- the total return swap.}
TRS-1's reference leg is defined as the NAV of a designated wallet, W-QIS, in a
virtual ledger. Valuing that leg requires no new machinery: it is the identical
fold, $\sum \mathrm{owned}\cdot P$ over W-QIS, that values any wallet. The
stamped NAV enters the true ledger as an observation, never a move
(Chapters~8 and~10), and the TRS contract reads it at each reset.
\emph{The reference leg of a swap is a NAV projection, folded like any other.}

\subsection{Dual valuation: two price vectors over one position set}

A wallet is valued for two purposes that demand different prices. Settlement,
margining, and reconciliation against a counterparty, a CCP, or a custodian
demand the price at which value actually changes hands --- the mark-to-market
price, $P^{\mathrm{mk}}$, which may differ by venue. Risk aggregation and PnL
attribution demand one internally consistent price across the book --- the
mark-to-mid price, $P^{\mathrm{md}}$. The ledger supplies both as projections
over one and the same composition.

\begin{principle}[One position set, two price vectors]
Mark-to-market and mark-to-mid are two price vectors applied to the single
composition produced by the fold. Position divergence is impossible: there is
one position record, so only two money figures --- $\mathrm{NAV}^{\mathrm{mk}}$
and $\mathrm{NAV}^{\mathrm{md}}$ --- can differ, never the positions
themselves.
\end{principle}

The classical source of internal reconciliation failure is that a settlement
system and a risk system each keep their own positions and drift apart. Here
they cannot: both read the same coordinates from the same log. The two numbers
that differ are honest reflections of two questions --- what would settle, and
what is the consistent risk mark --- and their difference is a basis, computed
on demand, never a break to be repaired.

\paragraph{Example --- three venue marks, one risk mid.}
A wallet holds $3{,}000$ units of $U$, sub-custodied across three venues:
$1{,}500$ at venue~A, $1{,}000$ at venue~B, $500$ at venue~C. Each venue
publishes its own official settlement close, used for that venue's margin and
reconciliation: $A = 100.20$, $B = 100.00$, $C = 99.60$. Risk marks the whole
holding at one mid, $100.00$.
\[
  \mathrm{NAV}^{\mathrm{mk}} = 1{,}500(100.20) + 1{,}000(100.00) + 500(99.60)
    = \USD{300,100},
\]
\[
  \mathrm{NAV}^{\mathrm{md}} = 3{,}000(100.00) = \USD{300,000}.
\]
The two money figures differ by \USD{100}, the cross-venue basis. The
$3{,}000$ is one number from one fold: each venue's sub-position reconciles
against its own statement under $P^{\mathrm{mk}}$, and the book aggregates for
risk under $P^{\mathrm{md}}$, but the position cannot diverge because there is
only one of it.

\subsection{The pricing-stack requirements}

The valuation discipline above forces a small set of requirements on the
pricing stack and the pricing data layer. They are derived here as consequences
of NAV being a projection; Chapter~16 restates them as the minima any
implementation must meet.

\begin{enumerate}[label=P\arabic*.,leftmargin=*]
\item \textbf{The ledger records pricing outputs and never runs a model.}
  Model outputs re-enter only as observations carrying sufficient provenance;
  reproducing a valuation from recorded prices is unconditional, while
  reproducing a model's number requires the model itself, whose retention is a
  governance matter outside scope.
\item \textbf{Valuation is state-parameterised, and the state comes from the
  ledger.} $P(u,\sigma(u))$ reads $\sigma(u)$ from UnitStatus and
  PositionState, never from a pricing-owned store --- the one-touch and the
  variance swap both turn on this.
\item \textbf{Valuations are reproducible from recorded inputs.} The
  composition is a fold of the log and the prices are identified on the record,
  so any party recomputes any NAV to the minor unit.
\item \textbf{Prices and quantities are never mixed across an event.} Each
  is taken in the same state of the instrument --- consistency of reference.
\item \textbf{Two price vectors are both projections over one composition.}
  Mark-to-market and mark-to-mid share the single position record; positions
  never diverge.
\item \textbf{Estimates are kept apart from entitlements}, and originals are
  preserved while adjusted and derived values are computed at read
  (Chapter~8).
\end{enumerate}

Each is a property that can be checked mechanically over generated products and
histories, and each is exercised by a thread episode in this chapter. Valuation
adds no store to the ledger: it is the log, folded and priced.
